\documentclass[11pt]{article}
\usepackage[T1]{fontenc}
\usepackage{textcomp}
\usepackage[margin=0.75in]{geometry}
\usepackage{amsmath,amssymb,amsthm}
\usepackage{array,booktabs}
\usepackage{hyperref}
\setlength{\parindent}{0pt}
\newtheorem{theorem}{Theorem}
\theoremstyle{definition}
\newtheorem{definition}{Definition}
\title{Flow Proof Artifact: Eq-subst-mul-right}
\author{Generated by \texttt{flow doc proof}}
\date{Flow Proof Book}
\begin{document}
\maketitle
Equal factors stay equal when multiplied on the right.\\[0.5em]
\textbf{Source.} leibniz — https://en.wikipedia.org/wiki/Substitution\_(logic)\\[0.5em]
\bigskip
\noindent{\large \textbf{Derived fact 1.} \textit{From a = b we deduce a * c = b * c} \hfill equality \cdot equality \cdot equal terms multiply the same on the right}\par
\smallskip\noindent\textit{Coordinate: equality \cdot equality \cdot equal terms multiply the same on the right}\par
\smallskip\noindent\textit{Source: peano/induction — Gries \& Schneider, Ch. 3}\par
\smallskip\noindent\textit{Built on: zero is the right annihilator, for multiplication on the natural numbers, successor on the right distributes, for multiplication on the natural numbers}\par
\medskip
\noindent\textbf{Goal.} From a = b we deduce a * c = b * c\par
\noindent$\displaystyle \forall a \in \mathbb{N} \forall b \in \mathbb{N} \forall c \in \mathbb{N}\quad a \cdot c = b \cdot c$\par
\medskip
\noindent
\begin{tabular}{@{} >{\raggedright\arraybackslash}p{0.06\textwidth} >{\raggedright\arraybackslash}p{0.47\textwidth} >{\raggedleft\arraybackslash}p{0.41\textwidth} @{}}
\textbf{\#} & \textbf{Proof} & \textbf{Mathematics} \\
\midrule
\textbf{1.} & We prove that equal terms multiply the same on the right for equality on equality. & \\[0.45em]
\textbf{2.} & We proceed by induction on a: first the base case, then the inductive step. & \\[0.45em]
\textbf{3.} & Consider the base case in which a  equals  b. & \\[0.45em]
\textbf{4.} & Consider the base case in which c is zero. & \\[0.45em]
\textbf{5.} & We invoke the definitional clause governing multiplication on the natural numbers: zero is the right annihilator, for multiplication on the natural numbers (instantiated for a). & \\[0.45em]
\textbf{6.} & We invoke the definitional clause governing multiplication on the natural numbers: zero is the right annihilator, for multiplication on the natural numbers (instantiated for b). & \\[0.45em]
\textbf{7.} & From step 4, step 5, and step 6, we can deduce that a times c equals b times c. This establishes the base case (see step 4, step 5, and step 6). Hence proven. & $\displaystyle a \cdot c = b \cdot c$ \\[0.45em]
\textbf{8.} & For the inductive step, suppose the claim holds for all smaller values. & \\[0.45em]
\textbf{9.} & Under the supposition in step 8, let k denote the predecessor of c. & \\[0.45em]
\textbf{10.} & Under the supposition in step 8, we cross the inductive boundary: assume the claim holds for a (the induction hypothesis). & $\displaystyle a \cdot k = b \cdot k$ \\[0.45em]
\textbf{11.} & We invoke the definitional clause governing multiplication on the natural numbers: successor on the right distributes, for multiplication on the natural numbers (instantiated for a, k). & \\[0.45em]
\textbf{12.} & We invoke the definitional clause governing multiplication on the natural numbers: successor on the right distributes, for multiplication on the natural numbers (instantiated for b, k). & \\[0.45em]
\textbf{13.} & From step 8, step 9, step 10, step 11, and step 12, this implies a times c equals b times c. Hence proven. & $\displaystyle a \cdot c = b \cdot c$ \\[0.45em]
\bottomrule
\end{tabular}
\label{thm:Eq-equality-equal_terms_multiply_the_same_on_the_right}
\par\medskip\hrule
\end{document}
